\section{Face phases and the last spectral cap}\label{lc:phases}
To use the signed real transfer at arbitrary face phases, we need
positivity of the interval continuant itself, rather than positivity
of its recurrence coefficients.

\begin{lemma}[Continuous real positivity of interval tails]\label{lc:tail-positive}
For \(0\le i\le j\le m\) and
\(0\le\delta\le1/2-1/(2N)\), \(K_{i,j}\) has nonnegative coefficients
in \(x,y\). Consequently
\[
 |K_{i,j}(X_0,Y_0;\delta)|
                 \le K_{i,j}(|X_0|,|Y_0|;\delta).
\]
\end{lemma}
\begin{proof}
Set \(U_r=K_{i,i+r}/\prod_{v=i+1}^{i+r}s_v^2\).
Then \(U_0=1,U_1=x+y\), and for \(r\ge2\), \(j=i+r\),
\[
 U_r=(x+y)U_{r-1}+\beta_jxyU_{r-2},\qquad
                       \beta_j=\frac{\alpha_{j-1}}{s_js_{j-1}}.
\]
The recurrence starts at its second step, so it never divides by \(s_0\).
Product-to-sum gives the row-independent identity
\[
 \alpha_{j-1}+s_js_{j-1}
             =4\cos^2(\pi\delta)-4\sin^2(\theta/2)\ge0.
\]
Thus \(1+\beta_j\ge0\), even if \(\beta_j<0\). Simultaneous induction
using
\[
 U_r-xU_{r-1}
   =y(U_{r-1}-xU_{r-2})+(1+\beta_j)xyU_{r-2}
\]
and \(U_0=1,U_1-xU_0=y\) proves positivity. Taking moduli in a
polynomial with nonnegative coefficients proves the last assertion.
\end{proof}

In \eqref{lc:left-summand}, take absolute values of the middle
monomials and of the scalar unused products. The lemma then permits
the positive outer moduli in the \emph{same signed} continuant.
Thus Proposition~\ref{lc:left-maximum} bounds every face phase
uniformly for
\[
                       -2\le\sigma\le b_N:=2\cos(\pi/N).
\]
No positivity of the full numerator, or comparison of its unused
products to their endpoint values, has been asserted.

For the small remaining interval, let \(p\) be any complex polynomial
of degree at most \(d\), bounded by \(M\) on \([-2,b_N]\).
The affine coordinate of \(\sigma=2\) relative to this interval is
\(z_N=(6-b_N)/(b_N+2)\).
For a degree-\(d\) polynomial \(r\), the function
\(w^{-d}r((w+w^{-1})/2)\) is holomorphic outside the unit disk and
at infinity. On its boundary it is bounded by \(\|r\|_{[-1,1]}\).
The maximum principle on the reciprocal disk gives
\[
 |r(z)|\le\|r\|_{[-1,1]}\exp(d\,\operatorname{arcosh}z),\qquad z\ge1.
\]
Complex coefficients cause no change. As \(z_N=1+O(N^{-2})\),
uniformly for \(N\ge3\),
\begin{equation}\label{lc:last-cap}
 \sup_{-2\le\sigma\le2}|p(\sigma)|
      \le M\exp(d\,\operatorname{arcosh}z_N)\le C M.
\end{equation}
For each fixed face-phase tuple the actual left polynomial has
spectral degree at most \(d\). Apply \eqref{lc:last-cap} before
taking the supremum of those phases. The constant is independent
of them. This proves the left assertion of Theorem~\ref{lc:main}.
Row reflection proves the right assertion. The central assertion
was proved in Section~\ref{lc:action}. Every compact estimate
respects the declared zero pattern, completing the theorem.

